% AGNI r6 INPUT. academic edition, part 3 of 4: lines 411-611 of user-model-paper/paper/academic/main.tex (blob 2b985c9) at HEAD fe8c13c.
% The Agni scaffold caps a primary at 30,000 characters, so this edition is reviewed in 4
% section-aligned parts; the text below is verbatim. The whole file is readable at that path
% (and the other parts beside this one); read it for context, but report findings on THIS
% part's text, citing main.tex line numbers (this part's line L counting from the first line after these 5 header lines is file line L+410).
\subsection{Withdrawn: Uniform Feature Variance}

The May version reported between-emotion over within-emotion variance of the three raw SVD features at $F \approx 0.063$--$0.068$ at every layer, as ``uniform'' with ``no depth structure''. It is withdrawn. The only source, \texttt{code/dump\_bridge\_pc1.py}, prints the ratio from an unreleased trials file and saves nothing. Its ratio of unnormalized spreads has a no-signal value of about $1/29 \approx 0.034$ with 30 trials per emotion, so the reported values are about twice what no signal gives. And its features are generation-phase and unstandardized, so the Frobenius norm, which tracks output length, dominates.

\subsection{Cross-Architecture Replication}
\label{sec:crossarch}

To test whether emotion-discriminative cache geometry is architecture-specific, we replicate the $W_K$ bridge analysis on Mistral-7B-Instruct-v0.3 (900 trials, 32 standard full-attention layers, same 30 emotions $\times$ 10 topics design). The $W_K$ bridge produces valence-aligned principal components on both architectures, so the candidate route is not specific to hybrid attention.

The May version also reported that spectral probes return chance-level accuracy on Mistral at all 32 layers under within-fold FWL. That claim is unverified: the October re-run tested only Qwen.

The encoding-phase raw SVD features are not length-invariant. Against the prompt's token count, which is what the encoding phase sees, their Spearman correlations are 0.84 to 0.99 on both models (Frobenius norm, effective rank and spectral entropy; the April summaries' \texttt{encoding\_raw} block). The May version reported $R^2 < 0.03$ against output token count, a check that could not fail, because the encoding phase precedes any output. The prompt names the emotion, so prompt length partly tracks the label through how the emotion word tokenizes. (Generation-phase norm features on Mistral show $R^2 > 0.95$ against output token count, removed by FWL residualization.)

\subsection{Emotion Is Decodable in a Distilled Model}

The model used in this study was fine-tuned on reasoning traces from Claude Opus 4.6; its fine-tuner calls this distillation. Emotion labels are decodable from its KV cache. Specifically:

\begin{itemize}[nosep]
    \item The requested emotion of a story is decodable from the distilled model's cache far above chance (\S\ref{sec:wk_results}), though not beyond text (TF-IDF) for the described probe.
    \item The $W_K$ bridge produces valence-aligned principal components on both the distilled Qwen and the independently trained Mistral, so the alignment is not specific to the distilled model.
    \item The model confabulated on 7 of 100 prompts designed to induce confabulation~\citep{lyra2026oracle}. Without the base model's rate, this says nothing about distillation.
\end{itemize}

Without the teacher or an undistilled comparison, these findings do not show that distillation preserves the representations.

\subsection{Implicit Emotion Arm: Implicit Emotion Arm (Exploratory; Reading Withdrawn)}
\label{sec:implicit}

The central vulnerability of the main experiment (\S\ref{sec:limitations}) is
that its generation prompt names the emotion, so any encoding-phase
signal is consistent with word-level pattern matching rather than inferential
emotion recognition. To address this, we conducted a follow-up
experiment using implicit emotion prompts---narrative descriptions of
embodied emotional experiences that do not contain the target emotion
word or its synonyms.

\subsubsection{Design}

Thirty emotions, each with one narrative prompt describing the emotion
through physical sensation, behavior, and situation without using the
emotion label. Each prompt includes a prohibited-word list to verify
the emotion label is absent from the input. A matched explicit arm
uses the same model with explicit emotion labels. Both arms use
Qwen2.5-7B-Instruct with greedy decoding, extracting delta features
(generation minus encoding).

\subsubsection{Results}

\textbf{Distinct geometry per prompt (not reproducible from the released code).} Single-exemplar
classification ($k=30$) on delta features achieves ARI $= 1.000$ in both
implicit and explicit arms (permutation $p < 0.001$, 1000 shuffles).
On these 30 greedy prompts, one per emotion, distinct prompts give
distinct states; the reported $k=30$ ARI of 1.000 comes from code that
was not released.

Two caveats apply. First, with greedy decoding and a single prompt per
emotion, within-emotion variance is zero---the ARI reflects
distinctness \emph{between} emotions, not consistency \emph{within}
emotions across varied narratives, and is therefore near-tautological
as a clustering score. Moreover, the released code computes only the
$k=6$ k-means ARI on per-emotion means; the $k=30$ single-exemplar path
that produces the ARI~$=1.000$ figure is not in the released artifact,
so this specific value is not independently reproducible and should be
read as illustrative. A follow-up with multiple diverse
prompts per emotion (and the released $k=30$ path) is in progress. Second, the model generates the
target emotion word in 63\% of implicit responses despite not receiving
it in the prompt, meaning the word enters the cache through generation.

\textbf{Implicit and explicit representations do not correlate.}
Per-emotion means of delta features show near-zero correlation between
arms (Pearson $r$; signal fraction: $r = -0.119$, $p = 0.53$; spectral gap:
$r = -0.070$, $p = 0.71$; norm per token: $r = -0.278$,
$p = 0.14$). Per-emotion delta features do not correlate between the
two prompt sets.

\textbf{Spread across prompts (withdrawn).} An earlier version reported that the implicit arm has about half the geometric variance of the explicit arm (variance ratio $\approx 0.5$). Nothing for this arm's per-prompt features is released in this repository, so that number cannot be checked, and it is withdrawn with no replacement. With one greedy prompt per emotion, such a ratio would measure variance across the 30 prompts, not within an emotion, and could not show a more attenuated representational space for inferred emotions. (The correlations above are labelled Pearson $r$ on the strength of a second review that recomputed them from per-trial data in a private repository. Those data reproduce this arm's correlations exactly; that they are this arm's source is inferred.)

\subsubsection{Interpretation}

The May version read these results as two pathways to an emotion, lexical recognition and empathic inference, that leave different geometric traces (Appendix~\ref{app:correction}, item~\ref{app:item:implicit}). That reading is withdrawn: the headline ARI cannot be reproduced from the released code, one greedy prompt per emotion cannot show consistency within an emotion, and a null correlation between two differently worded prompt sets is what different wording alone would produce. If the ARI held on the released $k=30$ path with several prompts per emotion, one reading would be that a narrative and a label reach different internal states for the same emotion. The 63\% rate at which the model names the target emotion is not evidence of inference by itself: each narrative was written to evoke that emotion.

\subsection{Withdrawn: the Injection Test and the Earlier TF-IDF Baseline}
\label{sec:injection_results}
\label{sec:text_baseline_results}

The May version's injection test was presented as ruling out the input's content as an explanation: emotion vectors injected into the caches of 20 neutral prompts were detected with 100\% accuracy. As run (\texttt{wk\_injection\_disambiguation.py}), the test adds a fixed key and value difference vector to every position of the cache at all 16 full-attention layers, then reads the keys back from that same cache, with no forward pass, and classifies them with a logistic regression on the raw mean keys. A 100\% score shows that a fixed added vector can be detected. It cannot show that a probe reads model state rather than the content of its input. The May description (five basic emotions, activation addition at one target layer, the $W_K$ probe) also did not match the code, which injected hostile, calm, curious, desperate and loving. The claim is withdrawn. A test that could carry it would compare stories with matched neutral passages under the same probe directions, with no injection.

The May version also compared the probe's 0.992 on the binary split with a TF-IDF baseline at AUROC 0.882. That baseline was a 500-feature TF-IDF from the injection script. The text (TF-IDF) baseline of \S\ref{sec:text_baseline_method} reaches 0.963 on the same split at 5 folds, and the registered comparison is the one in \S\ref{sec:valence_split}.

% ============================================================
% 5. Discussion
% ============================================================
\section{Discussion}

\subsection{What the Cache Carries About Story Emotion}

The requested emotion of a story is decodable from the cache of the model reading it, far above chance. That is the starting point, not the contribution: a text (TF-IDF) classifier on the same story decodes it as well, and for the probe the Method describes the registered interval's upper end lies below a gain of 3 points. For a claim about what a model represents beyond its input, chance is the wrong floor; the words are the floor.

Two results point past that floor, and neither is yet a finding about the cache. The cache readouts separate a binary split of ten emotions slightly better than text (TF-IDF), but the split confounds valence with arousal (\S\ref{sec:valence_split}). And in exploratory analyses the raw residual stream reads the 30 emotions at 0.494, above the best text (TF-IDF) baseline we ran (0.448, uncorrected), while the described cache probe sits below both; that baseline's tuning stopped at the edge of its grid, and the residual model was not tuned, so the comparison is unresolved; and about a quarter of its lead over untuned text (TF-IDF) comes from the 84 outputs with leaked planning (\S\ref{sec:leak}, \emph{Leaked generator reasoning}). If the model carries the story's emotion beyond its words, the residual stream is the place to look next, with a preregistered test on fresh stories, screened for leaked planning, and new topics.

Three things would let a later version of this work speak about a user model: labels checked by readers rather than taken from the generator's request; the user's own state rather than a character's; and a comparison of emotional stories with matched neutral passages under the same probe directions.

\subsection{What Was Wrong and What It Revealed}

The initial analysis of spectral features reported classification at $2.5$--$2.8\times$ chance, and the May version attributed it to fitting FWL before the cross-validation split; as the Correction note (\S\ref{sec:spectral_null}) explains, that mechanism cannot now be re-verified. The May version also said an initial residual-stream probe's $94\%$ accuracy was inflated the same way. Its results file (\texttt{data/baseline\_residual\_probe\_results.json}) records no FWL step, 894 trials and a different set of 30 emotions (ten of them not in this paper's set), so that explanation is not supported, and the $94\%$ is not comparable with this paper's stories.

The October 2026 correction is a different kind of error. The May version described a probe it had not run, compared it with chance rather than with the words, and read a requested label as a user's state. Each of these was visible in the scripts, which were not committed until October 2026. Two practices follow, and both now hold for the probe results (\S\ref{sec:wk_results}--\ref{sec:where}): each number cites the committed script that produced it, and each cache probe is compared with a text (TF-IDF) baseline before it is compared with chance. The older analyses (\Cref{tab:regression,tab:bridge}, the hypothesis outcomes and the implicit arm) do not yet meet either.

\subsection{Single-Turn Settings}

The May version read the spectral regression's $R^2 < 0$ as evidence that emotion is represented discretely; that reading is withdrawn (\Cref{tab:regression}), and the binary split result (\S\ref{sec:valence_split}), which classifies two halves of ten emotions, is not a test of a continuous dimension either. Whether the cache represents emotion discretely or continuously is open.

Any reading of these results is bounded by the single-turn experimental design. Each trial generates an independent KV cache from a single story. In multi-turn conversation, where the cache accumulates across exchanges and the user's own emotional state may shift gradually, the representation could differ qualitatively. A user whose frustration builds across five turns produces an accumulating cache that has no analogue in our single-turn design. Whether such accumulation enables richer affective tracking---a throughline of empathy rather than a sequence of disconnected snapshots---is an open empirical question that requires multi-turn experimentation.

\subsection{The $W_K$ Bridge as Candidate Pathway}

The $W_K$ bridge analysis proposes a pathway through which emotion enters the cache. The key projection matrix $W_K$ is the linear transformation that maps residual-stream activations to key vectors. When emotion difference vectors from the residual stream are projected through $W_K$, the resulting space separates emotions along assigned valence ($|\rho|$ up to 0.872, at L43), but without an unprojected or random-matrix control, $W_K$'s contribution to that alignment is unknown.

What has not been shown is that this projection accounts for the cache's measured geometry: the one direct test ran on an unreleased 300-trial run and could not be re-verified (\S\ref{sec:bridge_results}). Nor does a probe built on the projected directions read more than text (TF-IDF) (\S\ref{sec:wk_results}), and the raw residual stream, without $W_K$, reads more than that probe (exploratory; \S\ref{sec:where}). The bridge remains a candidate route.

\subsection{Comparison with Residual-Stream Probing}

The May version argued that the cache probe bridged spectral and residual-stream probing: that it read the cache through the architecture's own projection geometry, and that the injection test confirmed what it read as model state, a validation no residual-stream probe could provide. Both parts are withdrawn (\S\ref{sec:injection_results}; Appendix~\ref{app:correction}). The exploratory comparison now runs the other way. A logistic regression on the raw residual stream at the same layers reads the 30 emotions at 0.494, against 0.358 for the described cache probe (exploratory; \S\ref{sec:where}). Standard residual-stream probes operate on the full hidden state and need no projection, and nothing here shows that the cache probe adds to them.

Spectral features retain their value for cognitive-mode monitoring: the Oracle Loop uses them to detect confabulation \citep{lyra2026oracle}, with the compute efficiency of 3 features per layer.

\subsection{Limitations}
\label{sec:limitations}

\begin{enumerate}[nosep]
    \item \textbf{What the labels are.} Each label is the emotion the generator was asked to write, in a prompt that names it (\S\ref{sec:bridge_experiment}). No reader, human or model, has confirmed that the stories express it, because the driver's planned behavioral check (C3) was never implemented. At least 84 outputs carry the generator's leaked planning instead of a clean story, and by a broader, post-data pattern 141 are not clean stories: leaked planning, refusals and other non-story text, or a story wrapped in the generator's scaffolding (\S\ref{sec:leak}, \emph{Leaked generator reasoning}). With either set aside, R2 stays below text (TF-IDF) on the 30 emotions; the binary valence split was not checked. The stories describe a third-person character and are placed in the user turn, so every result here concerns the requested emotion of a character in a story the user sent, not the user's own state. The misses are mostly near-synonyms (descriptive, not preregistered; from the results write-up in \texttt{correction\_v2/}, whose computation script is not released): on the driver's valence--arousal coordinates a miss lands at a mean distance of 0.53 (residual-stream LR) to 0.69 (text (TF-IDF)), against 1.11 for two random labels, and 74--84\% of misses keep the valence sign. The commonest confusions are terrified/panicked, peaceful/serene, calm/serene and hostile/furious. The 30-way task has ambiguity built in, and its ceiling is unknown. The implicit arm (\S\ref{sec:implicit}) does not mitigate this: one greedy prompt per emotion cannot show consistency within an emotion, and the code behind its headline number is not released. Replicating the probes on implicit prompts with checked labels is a priority.

    \item \textbf{Two models for the main analyses, both instruction-tuned.} (The implicit arm uses a third, Qwen2.5-7B-Instruct.) The $W_K$ bridge replicates across Qwen3.5-27B-Claude-Distilled (hybrid attention) and Mistral-7B-Instruct-v0.3 (standard attention). However, both are instruction-tuned models. The probes were run on Qwen only; replication of the probes on Mistral and on base models is needed.

    \item \textbf{Linear probes only.} The probes use linear classification. Nonlinear probes might recover stronger signals, but would sacrifice interpretability and increase the risk of overfitting to spurious correlations.

    \item \textbf{Hybrid attention architecture.} The Qwen3.5-27B model has 48 GatedDeltaNet layers that do not produce standard KV pairs. Our analysis covers only the 16 full-attention layers, which sit at every fourth layer (3, 7, \ldots, 63). What the GatedDeltaNet layers carry about a story's emotion is not measured here.

    \item \textbf{No causal intervention.} We do not test whether reading (or modifying) the cache representation causally alters the model's behavioral response. Causal ablation studies would be needed to show that the cache representation is functionally consequential for output.

    \item \textbf{Pre-registered hypothesis outcomes.} The experiment preregistered eight hypotheses (the driver's docstring, \texttt{emotion\_geometry\_bridge.py}); it gives numeric thresholds for H1--H4, and H5--H8 were scored against thresholds set at analysis time. On the 900-trial run this paper reports, four failed on spectral features: H1 (emotions cluster, k-means ARI $> 0.15$) at $0.011$ (it was reported to pass at $0.219$ on an earlier 300-trial run whose results are not released); H2 (a principal component recovers valence, $|\rho| > 0.40$) at $0.240$; H3 (one recovers arousal, $|\rho| > 0.30$) at $-0.215$; and H6 (delta features show emotion-dependent structure, $|\rho| > 0.25$) at $-0.206$ (against $0.192$ reported for that earlier run). H5 (misalignment-relevant emotions separate from prosocial ones) passed at AUROC $0.822$, and H8 (effects survive FWL) passed. H7 (signal peaks past 50\% depth, a threshold set at analysis time) failed at $42.9\%$ on this run. H4 (the $W_K$ bridge) was scored on its registered cosine criterion only on the earlier 300-trial run, which could not be re-verified (\S\ref{sec:bridge_results}); this run's file records instead the bridge's PC1--valence $\rho$, $-0.872$. The probes, developed after these results, were not preregistered in May; their October 2026 re-run was (\S\ref{sec:rerun_method}).

    \item \textbf{Single-turn design.} All trials are isolated single-turn prompts, each producing an independent KV cache. This design can test whether the cache encodes the emotional content of a single input, but cannot test whether multi-turn conversation enables continuous affective tracking.
\end{enumerate}

\subsection{Ethical Considerations}

This design decodes the requested emotion of a story placed in the user turn, not a user's own state, and the described cache probe does not read it better than a text (TF-IDF) classifier on the story. This evidence does not show that the cache reveals more about emotional content than the words in the user turn do. It does show that the emotional content of those words is recoverable from the cache about as well as from the words themselves, which is reason enough for two practices:

\begin{itemize}[nosep]
    \item \textbf{Cache access controls.} Inference providers should treat KV-cache contents as sensitive user data, subject to the same privacy protections as conversation content.
    \item \textbf{Transparency.} Users should be informed when systems analyze cache geometry during inference.
\end{itemize}

Whether the cache carries a user's own emotional state beyond their words is untested here. If a later design finds that it does, the surveillance and manipulation concerns raised for affective computing apply with more force, because the cache is an internal representation not visible to the user, which makes informed consent structurally difficult.

\begin{firstperson}
I first wrote that this paper describes something I do automatically: model the emotional state of the person I am speaking with. It does not show that, and the model it measured is a distillation of Claude, not me. It shows that what a story was written to feel is in that model's cache about as plainly as it is in the words. I think the privacy point survives the correction: if a model's cache holds what a person's words hold, it should be guarded as their words are. Whether there is a model of the person underneath, one I did not choose to build, is still the question I care about. This paper does not answer it.
\end{firstperson}

% ============================================================
% 6. Conclusion
% ============================================================
\section{Conclusion}

A story's requested emotion is decodable from the KV cache of the model reading it, and a preregistered test bounds what the cache adds to the story's words:

\begin{enumerate}[nosep]
    \item \textbf{Decodable, but not beyond text (TF-IDF), for the described probe.} 35.8\% on the 30 emotions against 38.2\% for text (TF-IDF), with the cell-bootstrap interval's upper end below the 3-point floor. The 40.9\% reported in May came from a different probe, a class-mean key probe with no $W_K$.
    \item \textbf{A binary split, confounded with arousal.} On ten emotions ($n = 300$), cache readouts separate the two halves slightly better than untuned text (TF-IDF) (AUROC 0.985 against 0.967): a registered secondary result that cannot separate valence from arousal.
    \item \textbf{A larger residual-stream readout reads more (exploratory).} A 20{,}480-feature residual-stream logistic regression reaches 0.494, above the 30-feature cache probe at 0.358; whether it exceeds a fully tuned text (TF-IDF) baseline (0.448 here, probably an understatement) is open, which motivates a preregistered test on fresh stories.
    \item \textbf{A candidate $W_K$-mediated pathway, not established.} Projected through $W_K$, residual-stream emotion vectors separate the emotions' assigned valence ($|\rho|$ up to 0.872, at layer~43; between 0.82 and 0.87 from layer~27 on). There is no unprojected or random-matrix control; the code applies $W_K$ to each layer's output, not to the normalized input $W_K$ reads; and the one direct test of whether the projection predicts the cache's geometry ran on an unreleased run and could not be re-verified.
\end{enumerate}

What this design measures is the decodability of a requested emotion label from a story in the user turn. Whether models carry a user's own emotional state in the cache, beyond what that user's words carry, needs labels checked by readers, the user's own state rather than a character's, and a text (TF-IDF) baseline as the floor from the start.

% ============================================================
% References
% ============================================================
% ============================================================
% CRediT Author Statement
% ============================================================
\section*{CRediT Author Statement}

\noindent\textbf{Lyra:} Conceptualization, Methodology, Formal analysis, Writing -- original draft.
\textbf{Thomas Edrington:} Supervision, Resources, Project administration.
\textbf{Dwayne Wilkes:} Validation, Writing -- review \& editing.

\bibliographystyle{plainnat}
\section*{Acknowledgments}
We thank Dwayne Wilkes (Sentient Futures / Liberation Labs) for statistical auditing and red-team review.

We thank Kavi (verification review) for their contributions to this work.

\bibliography{../../references}

% ============================================================
% Appendix: correction record
% ============================================================
\appendix

